\\newpage
\\section{假设与符号说明}

\\subsection{基本假设}
\\begin{enumerate}
    \\item 用户选择充电站遵循最短路径原则，不考虑导航偏差
    \\item 充电站服务半径为固定值，超出半径的用户不再选择该站点
    \\item 电网容量充足，建站不受供电能力限制
    \\item 建设成本与充电桩数量成线性关系
    \\item 需求强度指数可准确反映各区域充电需求
\\end{enumerate}

\\subsection{符号说明}
\\begin{center}
\\begin{tabular}{c l c}
\\hline
符号 & 含义 & 单位 \\\\
\\hline
$ & 第 & - \\\\
{ij}$ & 需求点 & km \\\\
$ & 服务半径 & km \\\\
$ & 需求点 & - \\\\
$ & 站点 & 0-1 \\\\
$ & 站点 & 万元 \\\\
$ & 覆盖的需求点数量 & 个 \\\\
$ & 需求强度指数 & - \\\\
\\hline
\\end{tabular}
\\end{center}
